\documentclass[11pt]{article}
\usepackage{fontspec}
\setmainfont{Times New Roman}
\setsansfont{Arial}
\setmonofont{Consolas}
\usepackage{amsmath,amssymb}
\usepackage{geometry}
\usepackage{microtype}
\geometry{a4paper,margin=3cm}
\setlength{\parindent}{1.25em}
\setlength{\parskip}{0pt}
\emergencystretch=8em
\allowdisplaybreaks
\newcommand{\pair}[2]{\left(#1\middle|#2\right)}
\newcommand{\eps}{\varepsilon}

\begin{document}

\section*{\S\ 6. Експлицитно представјенье инвариантных творјень.\\Процесы свијаньа и диференциације.}

Резултаты предходных параграфов дозвольајут нам довести до конца теорему о симболичној представимости, изречену в \S\ 2 (прву фундаменталну теорему), додаваньем експлицитној представимости; именно покажемо:

\emph{Теорема IV:} ``Все инвариантне творјеньа либоволных основных форм можно експлицитно представити матричными продуктами; то јест в закльучне изразы, осим рядов пременных $p_\rho$, входе само ряды $S^\sigma,\ldots,T^\tau$ првотных форм или ряды $P,Q,R$, изведене из ньих субституцијеју (9.) или (11.).''

За доказ предпоставимо, же инвариантно творјенье, осим рядов пременных, може зависети од рядов $S^\sigma,\ldots,T^\tau$; и то нехај ти ряды сут достаточно далеко разложене на поједине ряды:
$S^\sigma=S^{\sigma_1}\cdots S^{\sigma_i},\ldots,T^\tau=T^{\tau_1}\cdots T^{\tau_k}$, абы было можно представити јих детерминантами. Ряды нехај сут уреджене в самом по себе либоволном, но точно установјеном порядку $S^\sigma,\ldots,T^\tau$. Изберемо таковы агрегат детерминантов, кторы зависи од првого целого реда $S^\sigma=S^{\sigma_1}\cdots S^{\sigma_i}$, а не само од појединых јего частных рядов. Та зависност јест по \S\ 3--\S\ 5 последком обчих идентичностиј. Ако тепер разложимо ряд пременных $p_\rho$ в ряды $y,v$, или једен из позднејших рядов $T$ в поједине ряды $t$, односно $\tau$, тогда најобчејше идентичности сут составјене из (20.), (21.). Но те последнье идентичности всегда воде к изразу, кторы експлицитно содержи ряд $S^\sigma=S^{\sigma_1}\cdots S^{\sigma_i}$, то јест леву страну (20.), (21.); чрез (19.) се имено уверујемо, же тек полна десна сума -- а не јеј част, ктора одговарја члену (20а.) -- имаје својство зависети само од $S^\sigma$. При продолженьу поступка все ряды, кторе једин раз входе експлицитно, оставајут експлицитно сохраньене; применьенье идентичностиј може једино евентуално требовати сојединьенье неколико рядов в једен, то јест субституцију (9.). Ако наконец треба привести к експлицитној форме само једен ряд $T^\tau$, можне сут два случаје. Може јешче наступити свободны ряд пременных, то јест ряд не связаны с другими рядами субституцијеју (9.); ньеговым разложеньем в частне ряды $y$ или $v$ можно довести поступак до конца, и наставајут, како в (31.), полары једној или неколико форм, кторе експлицитно содержет все ряды. Ако пак свободны ряд пременных више не наступаје, тогда по начину ньих настаньа познајемо в целости изразов, кторе содержет поједине ряды $t$ или $\tau$, но зависят од реда $T^\tau$, члены разкладној идентичности; те пак по \S\ 5 при применьеньу субституције (11.) всегда допушчајут експлицитно представјенье во всех рядах. Тым јест горньа теорема обче доказана. Треба јешче заметити, же когда наступајут само два ряды симболов, субституција (11.) никогда не наступаје; бо -- из-за $\sigma,n-\sigma,\tau,n-\tau<n$ -- ряды пременных могут не входити само тогда, когда оба ряды симболов сут сјединьене в једном детерминанту, теды наступајут експлицитно; но как скоро входе ряды пременных, (34.) и (33.) показујут, же субституција (11.) јест непотребна.

Из управо реченого следује, же за произведенье всех инвариантных творјень достаточно јест сојединьати ряды симболов, с додаваньем рядов пременных, в все можне матричне продукты. Најобчејши матричны продукт имаје форму:
\begin{equation}
 \pair{P^{\mu_1}\cdots P^{\mu_i}}{Q_{\nu_1}\cdots Q_{\nu_k}},
 \qquad \sum\mu=\sum\nu,
\tag{40}
\end{equation}
кде $P$ и $Q$ могут быти ряды првотных форм, или могут настати из првотных рядов чрез последовност субституциј (9.) или (11.), или наконец быти ряды пременных. Изразы (40.) пак можно произвести поступным сојединьеньем каждах двох рядов симболов в матричны продукт, с помочју рядов пременных; бо из
\[
 \pair{q_{n-\mu-\lambda}P^\mu}{Q_{\nu_1}p_{n-\mu-\nu_1-\lambda}}
 =\pair{R_\lambda Q_{\nu_1}}{p_{n-\mu-\nu_1-\lambda}}
\]
сојединьеньем с редом $Q_{\nu_2}$ достајемо израз
\[
 \pair{[R_\lambda Q_{\nu_1}]^{n-\lambda-\nu_1}}
 {Q_{\nu_2}p_{n-\mu-\nu_1-\nu_2-\lambda}}
 =\pair{q_{n-\mu-\lambda}P^\mu}
 {Q_{\nu_1}Q_{\nu_2}p_{n-\mu-\nu_1-\nu_2-\lambda}},
\]
и продолженьем поступка достајемо израз (40.). Ако при том $P$ и $Q$ не сут ряды првотных форм, тогда (9.) и (11.) в связи с управо показаным доказујут, же они такоже настали чрез последовност сојединьень каждах двох рядов симболов. Из тога следује:

\emph{Теорема V:} Процес произведеньа всех инвариантных творјень, процес свијаньа, состаји се из поступного сојединьеньа каждах двох рядов симболов в матричны продукт с помочју рядов пременных.

Из двох рядов симболов $S^\sigma,T^\tau$, кде $\sigma\ge\tau$, можно утворити следујуче матричне продукты:
\begin{align}
&\pair{q_{n-\sigma-\lambda}S^\sigma}{T_{n-\tau}p_{\tau-\lambda}},
&&(\lambda=1,2,\ldots,\min(\tau,n-\sigma)),\tag{41}\\
&\pair{q_{n-\tau-\mu}T^\tau}{S_{n-\sigma}p_{\sigma-\mu}},
&&(\mu=\sigma-\tau+\lambda),\tag{42}\\
&\pair{q_{n-\tau-\nu}T^\tau}{S_{n-\sigma}p_{\sigma-\nu}},
&&(\nu=1,2,\ldots,\sigma-\tau).
\tag{43}
\end{align}
Из (31.) пак за (42.) и (43.) следујут вредности:
\begin{align}
&\pair{q_{n-\sigma-\lambda}S^\sigma}{T_{n-\tau}p_{\tau-\lambda}}
 +\sum_{\alpha}\Pi
 \pair{q_{n-\sigma-\lambda-\alpha}S^\sigma}
 {T_{n-\tau}p_{\tau-\lambda-\alpha}},\tag{42a}\\
&\Pi(q_{n-\sigma}S^\sigma)(T_{n-\tau}p_\tau)
 +\sum_{\beta}\Pi
 \pair{q_{n-\sigma-\beta}S^\sigma}{T_{n-\tau}p_{\tau-\beta}}.
\tag{43a}
\end{align}
Теды формы (42.) можно линеарно изразити чрез (41.) и полары од (41.), а формы (43.) чрез полары основној формы и форм (41.). Такоже из (31.) следује линеарна зависност форм (41.) од (42.) и поларов од (42.). По дефиницији Клебша\footnote{Clebsch, а. а. О. \S\ 7.} формы (42.) сут еквивалентне формам (41.), доколе (43.) врача к основној форме. Производечи процес јест теды равноправно даны чрез (41.) или (42.); процес (41.), простејши с обзиром на число $\lambda$, будемо определьати како процес свијаньа. Число $\lambda$, дефект свијаньа (ср. \S\ 2), даје число рядов, кторых недостаје до детерминантного продукта, и то в оном матричном продукту, кторы при даном свијаньу содержи максимално число рядов. Понеже $\lambda$ бежи само до меншеј из двох вредностиј $\tau,n-\sigma$, кде $\sigma\ge\tau$, јест:
\begin{equation}
 \lambda\le\frac n2,\ n\text{ парно};\qquad
 \lambda\le\frac{n-1}{2},\ n\text{ непарно}.
\tag{44}
\end{equation}
Сведенье (42.) и (43.) к (41.) оставаје важече и тогда, когда чрез дальше свијанье ряды $p,q$ прешли в ряды симболов; бо по (36.) достајемо:
\[
 \pair{A^{n-\tau-\varkappa}T^\tau}{S_{n-\sigma}B_{\sigma-\varkappa}}
 =\sum_{\alpha=\max(0,\varkappa-\sigma+\tau)}^{\min(\tau,n-\sigma)}
 \eps\pair{R^{\tau-\alpha}B^{n-\sigma+\varkappa}}
 {A_{\tau+\varkappa}R_{n-\sigma-\alpha}},
\]
кде ряды $R$ настали из $S$ и $T$ по (33.). Теды те формы, кторе содержет само ряды симболов, добивајут се свијаньем $\pair{q_\tau A^{n-\tau-\varkappa}}{B_{\sigma-\varkappa}p_{n-\sigma}}$ односно $\pair{q_{\tau-\alpha}B^{n-\sigma+\varkappa}}{A_{\tau+\varkappa}p_{n-\sigma-\alpha}}$ -- в зависности од тога, чи $\sigma-\tau\gtreqless 2\varkappa$ -- с основној формоју и формами (41.).

Из симболичного процеса свијаньа в познаты способ добивајут се инвариантне процесы диференциације, когда само ряды симболов $S^\sigma,T_{n-\tau}$ заменимо рядами диференциалных симболов $\frac{\partial}{\partial p_\sigma}$, $\frac{\partial}{\partial q_{n-\tau}}$, при чем, како познато, всакому продукту $\frac{\partial}{\partial p_{i_1\ldots i_\sigma}}\frac{\partial}{\partial q_{j_1\ldots j_{n-\tau}}}$ одговарја реално значенье $\frac{\partial^2}{\partial p_{i_1\ldots i_\sigma}\partial q_{j_1\ldots j_{n-\tau}}}$. Понеже ряды $\frac{\partial}{\partial p_\sigma}$, $\frac{\partial}{\partial q_{n-\tau}}$ сут когредиентне редам $S^\sigma,T_{n-\tau}$, они изполньајут исте идентичности како те, теды специално такоже идентичности меджу (42.) и (42а.), (43.) и (43а.). Составјајучи резултат, достајемо:

\emph{Теорема VI:} Все инвариантне творјеньа либоволных основных форм наставајут из ньих чрез повторно применьенье симболичного процеса свијаньа (41.) или такоже инвариантных процесов диференциације:
\begin{equation}
 \left(q_{n-\sigma-\lambda}\frac{\partial}{\partial p_\sigma}
 \middle|\frac{\partial}{\partial q_{n-\tau}}p_{\tau-\lambda}\right),
 \qquad(\sigma\ge\tau,\ \lambda=1,2,\ldots,\min(\tau,n-\sigma)).
\tag{45}
\end{equation}
Треба јешче заметити, же при всаком свијаньи (41.) вкупна вага формы, то јест сума ваг појединых рядов пременных (ср. \S\ 1), уменшаје се за позитивно число $2(\lambda^2+\lambda(\sigma-\tau))$. Из тога, независно од обчего доказа конечности, следује конечност реда форм, то јест целости инвариантных творјень данеј основној формы, линеарных в коефициентах.\footnote{Ср. Clebsch, а. а. О. \S\ 11 и 12. Конечност система елементарных ковариантов.}

Далье треба заметити, же теорема VI такоже важи за формы, кторе не изполньајут условја (12.). Бо всакы фактор из (3.), $\pair{S^\rho}{p_\rho}^{m_\rho}$, можно заменити продуктом $m_\rho$ линеарных факторов $\pair{A^\rho}{p_\rho}\pair{B^\rho}{p_\rho}\cdots\pair{M^\rho}{p_\rho}$; тым инвариантне творјеньа преходят в творјеньа система линеарных форм, кторе тек позднеје треба поставити равными једна другој. За всаку поједину линеарну форму пак условја (12.), кторе при уведеньу диференциалных симболов содержет само частне дериваты 2-го реда, всегда сут изполньене.

\end{document}
